\scanpage{20}\mspage{\unsure{19|17}}
% @ first three lines
pour arriver à la notion de lieu, et d'interprétation des données en termes de lieux.

% @ terminology
\unsure{Pt|P\textsuperscript{t}} de vue \unsure{terminologie|terminologique}, au lieu de ``figure élémentaire'', je préfère dire maintenant ``\uline{multistrate}'', plus parlant.

Donc une figure quelconque est un ens de multistrates mutuellement compatibles, fermé pour $\leq$, et loc. fini.

% @ definition of lieu
\deleted{\unsure{F|E}} On appelle \deleted{\unsure{multi}} \uline{lieu} de la contrée, une multistrate qui est minimale pour $\ll$, i.e. qui n'a d'autre raffinement qu'elle-\unsure{\={m}|m}. \deleted{C'est donc} \deleted{\unsure{point|\c{c}a} \unsure{d'un pt}} \deleted{C'est donc} \deleted{Dans le} (Dans les exemples-types, c'est un ens de parties de $\deleted{\unsure{A}}X$, réduits à une seule partie, qui est \unsure{ponctuelle} -- \unsure{donc|dans} les lieux sont les p\textsuperscript{ts} de $X$ \unsure{..|.})

% @ multistrates minimales
Ce sont donc aussi des multistrates minimales pour $\leq$, \deleted{\unsure{donc}} \inserted{à fortiori} des ``\unsure{variétés|vacuités}'' \deleted{\unsure{connexes} \unsure{non} \unsure{vides})''} -- ce sont les non-variétés \unsure{0-connexes|o-connexes} qui n'ont pas de raffinement autre qu'elle-\={m}.

% @ last lines and inclusion
Soit $\unsure{L|\mathcal{L}}$ l'ensemble des lieux, \deleted{Pour tout} donc
\[
\begin{aligned}
L\subset M\quad (={}&F_{el}=\text{ensemble des}\\
&\hphantom{F_{el}={}}\text{multistrates}\,)
\end{aligned}
\]
